% GENERATED FILE. Edit canonical lessons, then run npm run generate:books.
% canonical-source-hash: fnv1a64:7c513d048d4b3ff2
% canonical-lessons: TA-C424-catukkam, TA-C424-cilai, TA-C424-ninaivuccinnam, TA-C424-malikai, TA-C424-tutarakam

\chapter{Statues and Monuments}
\label{ch:ta-statues-and-monuments}
\hlchaptermodality{424}

\begin{chapteropening}
\textbf{By the end of this chapter:} \emph{I can talk about a town square, statues, monuments, a mansion and an embassy.}

Complete the last lesson of chapter 424: I can talk about a town square, statues, monuments, a mansion and an embassy.
\end{chapteropening}

\section[\texorpdfstring{\ta{சதுக்கம்} (catukkam)}{சதுக்கம் (catukkam)}]{\ta{சதுக்கம்} (catukkam) — a town square}

\label{lesson:TA-C424-catukkam}



\begin{quote}
Before the new one: say the Tamil for a locality, then the Tamil for a centre.
\end{quote}

\subsection*{You'll want to know: \ta{சதுக்கம்}}
\textbf{\ta{சதுக்கம்}} — \emph{catukkam} — \textquotedblleft{}a town square\textquotedblright{}.

\begin{quote}
\ta{சதுக்கத்தில்} \ta{நிறைய} \ta{மக்கள்} \ta{கூடியிருந்தார்கள்}.
\end{quote}

\begin{quote}
\emph{catukkattil niṟaiya makkaḷ kūṭiyiruntārkaḷ.}
\end{quote}

\begin{quote}
Many people had gathered in the square.
\end{quote}

Before a case the -\ta{அம்} becomes -\ta{அத்து}-: \ta{சதுக்கத்தில்}.

\begin{cousinweb}
\textbf{\ta{சதுக்கம்}} (\emph{catukkam}) is borrowed from Prakrit \emph{catukka}, \textquotedblleft{}a set of four, a crossroads, a village square\textquotedblright{}, from Sanskrit \emph{catuṣka} (CDIAL 4629).
\end{cousinweb}

\begin{grammarlens}[title={What you've built}]
One more word for getting around.
\end{grammarlens}

\subsection*{Your turn}
\begin{itemize}
\raggedright
  \item \emph{Say it:} \emph{catukkam}
  \item \emph{Say it:} \emph{catukkam}, once more
  \item \emph{Say it:} \emph{catukkattil niṟaiya makkaḷ kūṭiyiruntārkaḷ.} — Many people had gathered in the square.
  \item \emph{Recall:} say the Tamil for a locality, then the Tamil for a centre, then say \emph{catukkam} again
\end{itemize}

\begin{tcolorbox}[breakable,colback=teal!4,colframe=teal!35!black,title={Before you move on}]
What does \ta{சதுக்கம்} mean? (\textquotedblleft{}A town square\textquotedblright{}.) Say it once more.
\end{tcolorbox}
\section[\texorpdfstring{\ta{சிலை} (cilai)}{சிலை (cilai)}]{\ta{சிலை} (cilai) — a statue}

\label{lesson:TA-C424-cilai}



\begin{quote}
Before the new one: say the Tamil for a centre, then the Tamil for a town square.
\end{quote}

\subsection*{You'll want to know: \ta{சிலை}}
\textbf{\ta{சிலை}} — \emph{cilai} — \textquotedblleft{}a statue\textquotedblright{}.

\begin{quote}
\ta{சதுக்கத்தின்} \ta{நடுவில்} \ta{ஒரு} \ta{பெரிய} \ta{சிலை} \ta{இருக்கிறது}.
\end{quote}

\begin{quote}
\emph{catukkattiṉ naṭuvil oru periya cilai irukkiṟatu.}
\end{quote}

\begin{quote}
There is a big statue in the middle of the square.
\end{quote}

The plural is \ta{சிலைகள்}.

\begin{cousinweb}
\textbf{\ta{சிலை}} (\emph{cilai}) is borrowed from Sanskrit \emph{śilā}, \textquotedblleft{}a rock, a stone\textquotedblright{} (CDIAL 12459): a figure carved in stone.
\end{cousinweb}

\begin{grammarlens}[title={What you've built}]
One more word for getting around.
\end{grammarlens}

\subsection*{Your turn}
\begin{itemize}
\raggedright
  \item \emph{Say it:} \emph{cilai}
  \item \emph{Say it:} \emph{cilai}, once more
  \item \emph{Say it:} \emph{catukkattiṉ naṭuvil oru periya cilai irukkiṟatu.} — There is a big statue in the middle of the square.
  \item \emph{Recall:} say the Tamil for a centre, then the Tamil for a town square, then say \emph{cilai} again
\end{itemize}

\begin{tcolorbox}[breakable,colback=teal!4,colframe=teal!35!black,title={Before you move on}]
What does \ta{சிலை} mean? (\textquotedblleft{}A statue\textquotedblright{}.) Say it once more.
\end{tcolorbox}
\section[\texorpdfstring{\ta{நினைவுச்சின்னம்} (niṉaivucciṉṉam)}{நினைவுச்சின்னம் (niṉaivucciṉṉam)}]{\ta{நினைவுச்சின்னம்} (niṉaivucciṉṉam) — a monument; a memorial}

\label{lesson:TA-C424-ninaivuccinnam}



\begin{quote}
Before the new one: say the Tamil for a town square, then the Tamil for a statue.
\end{quote}

\subsection*{You'll want to know: \ta{நினைவுச்சின்னம்}}
\textbf{\ta{நினைவுச்சின்னம்}} — \emph{niṉaivucciṉṉam} — \textquotedblleft{}a monument; a memorial\textquotedblright{}.

\begin{quote}
\ta{இந்த} \ta{நினைவுச்சின்னம்} \ta{நூறு} \ta{ஆண்டுகள்} \ta{பழையது}.
\end{quote}

\begin{quote}
\emph{inta niṉaivucciṉṉam nūṟu āṇṭukaḷ paḻaiyatu.}
\end{quote}

\begin{quote}
This monument is a hundred years old.
\end{quote}

Two words joined, with the \ta{ச்} doubled: \ta{நினைவு} + \ta{சின்னம்}.

\begin{cousinweb}
\textbf{\ta{நினைவுச்சின்னம்}} (\emph{niṉaivucciṉṉam}) is built from \emph{niṉaivu}, \textquotedblleft{}remembrance\textquotedblright{} (DEDR 3683), and \ta{சின்னம்}, \textquotedblleft{}a sign\textquotedblright{}, borrowed from Sanskrit \emph{cihna}, \textquotedblleft{}a mark\textquotedblright{} (CDIAL 4833).
\end{cousinweb}

\begin{grammarlens}[title={What you've built}]
One more word for getting around.
\end{grammarlens}

\subsection*{Your turn}
\begin{itemize}
\raggedright
  \item \emph{Say it:} \emph{niṉaivucciṉṉam}
  \item \emph{Say it:} \emph{niṉaivucciṉṉam}, once more
  \item \emph{Say it:} \emph{inta niṉaivucciṉṉam nūṟu āṇṭukaḷ paḻaiyatu.} — This monument is a hundred years old.
  \item \emph{Recall:} say the Tamil for a town square, then the Tamil for a statue, then say \emph{niṉaivucciṉṉam} again
\end{itemize}

\begin{tcolorbox}[breakable,colback=teal!4,colframe=teal!35!black,title={Before you move on}]
What does \ta{நினைவுச்சின்னம்} mean? (\textquotedblleft{}A monument; a memorial\textquotedblright{}.) Say it once more.
\end{tcolorbox}
\section[\texorpdfstring{\ta{மாளிகை} (māḷikai)}{மாளிகை (māḷikai)}]{\ta{மாளிகை} (māḷikai) — a mansion; a palace}

\label{lesson:TA-C424-malikai}



\begin{quote}
Before the new one: say the Tamil for a statue, then the Tamil for a monument.
\end{quote}

\subsection*{You'll want to know: \ta{மாளிகை}}
\textbf{\ta{மாளிகை}} — \emph{māḷikai} — \textquotedblleft{}a mansion; a palace\textquotedblright{}.

\begin{quote}
\ta{அந்தப்} \ta{பழைய} \ta{மாளிகை} \ta{இப்போது} \ta{அருங்காட்சியகம்}.
\end{quote}

\begin{quote}
\emph{antap paḻaiya māḷikai ippōtu aruṅkāṭciyakam.}
\end{quote}

\begin{quote}
That old mansion is now a museum.
\end{quote}

The plural is \ta{மாளிகைகள்}.

\begin{cousinweb}
\textbf{\ta{மாளிகை}} (\emph{māḷikai}) is borrowed from Sanskrit \emph{mālikā}, which the Sanskrit lexicographers use for a tall house with an upper storey (Monier-Williams).
\end{cousinweb}

\begin{grammarlens}[title={What you've built}]
One more word for getting around.
\end{grammarlens}

\subsection*{Your turn}
\begin{itemize}
\raggedright
  \item \emph{Say it:} \emph{māḷikai}
  \item \emph{Say it:} \emph{māḷikai}, once more
  \item \emph{Say it:} \emph{antap paḻaiya māḷikai ippōtu aruṅkāṭciyakam.} — That old mansion is now a museum.
  \item \emph{Recall:} say the Tamil for a statue, then the Tamil for a monument, then say \emph{māḷikai} again
\end{itemize}

\begin{tcolorbox}[breakable,colback=teal!4,colframe=teal!35!black,title={Before you move on}]
What does \ta{மாளிகை} mean? (\textquotedblleft{}A mansion; a palace\textquotedblright{}.) Say it once more.
\end{tcolorbox}
\section[\texorpdfstring{\ta{தூதரகம்} (tūtarakam)}{தூதரகம் (tūtarakam)}]{\ta{தூதரகம்} (tūtarakam) — an embassy}

\label{lesson:TA-C424-tutarakam}



\begin{quote}
Before the new one: say the Tamil for a monument, then the Tamil for a mansion.
\end{quote}

\subsection*{You'll want to know: \ta{தூதரகம்}}
\textbf{\ta{தூதரகம்}} — \emph{tūtarakam} — \textquotedblleft{}an embassy\textquotedblright{}.

\begin{quote}
\ta{விசாவுக்காக} \ta{நாங்கள்} \ta{தூதரகத்துக்குப்} \ta{போனோம்}.
\end{quote}

\begin{quote}
\emph{vicāvukkāka nāṅkaḷ tūtarakattukkup pōṉōm.}
\end{quote}

\begin{quote}
We went to the embassy for a visa.
\end{quote}

\ta{தூதர்}, an envoy, + \ta{அகம்}, a house.

\begin{cousinweb}
\textbf{\ta{தூதரகம்}} (\emph{tūtarakam}) is built from \ta{தூதர்}, \textquotedblleft{}an envoy\textquotedblright{}, borrowed from Sanskrit \emph{dūta}, \textquotedblleft{}a messenger\textquotedblright{} (CDIAL 6488), and \ta{அகம்}, \textquotedblleft{}the inside, a house\textquotedblright{}.
\end{cousinweb}

\begin{grammarlens}[title={What you've built}]
One more word for getting around.
\end{grammarlens}

\subsection*{Your turn}
\begin{itemize}
\raggedright
  \item \emph{Say it:} \emph{tūtarakam}
  \item \emph{Say it:} \emph{tūtarakam}, once more
  \item \emph{Say it:} \emph{vicāvukkāka nāṅkaḷ tūtarakattukkup pōṉōm.} — We went to the embassy for a visa.
  \item \emph{Recall:} say the Tamil for a monument, then the Tamil for a mansion, then say \emph{tūtarakam} again
\end{itemize}

\begin{tcolorbox}[breakable,colback=teal!4,colframe=teal!35!black,title={Before you move on}]
What does \ta{தூதரகம்} mean? (\textquotedblleft{}An embassy\textquotedblright{}.) Say it once more.
\end{tcolorbox}
